\documentclass[preprint,12pt,authoryear]{elsarticle}
%DIF LATEXDIFF DIFFERENCE FILE



%% Use the option review to obtain double line spacing
%% \documentclass[authoryear,preprint,review,12pt]{elsarticle}

%DIF 6-10d6
%DIF < %% For including figures, graphicx.sty has been loaded in
%DIF < %% elsarticle.cls. If you prefer to use the old commands
%DIF < %% please give \usepackage{epsfig}
%DIF < 
%DIF < %% The amssymb package provides various useful mathematical symbols
%DIF -------
\usepackage{amssymb}
\linespread{1.2}
\usepackage{amsmath}
\usepackage{lineno}
\usepackage{graphicx}
\usepackage[colorlinks=true, allcolors=blue]{hyperref}
%DIF 17c12
%DIF < \usepackage[a4paper, total={6.in, 9in}]{geometry}
%DIF -------
\usepackage[a4paper, total={6in, 9in}]{geometry} %DIF > 
%DIF -------
% \linenumbers
\journal{Estuaries and Coasts}
%DIF PREAMBLE EXTENSION ADDED BY LATEXDIFF
%DIF UNDERLINE PREAMBLE %DIF PREAMBLE
\RequirePackage[normalem]{ulem} %DIF PREAMBLE
\RequirePackage{color}\definecolor{RED}{rgb}{1,0,0}\definecolor{BLUE}{rgb}{0,0,1} %DIF PREAMBLE
\providecommand{\DIFaddtex}[1]{{\protect\color{blue}\uwave{#1}}} %DIF PREAMBLE
\providecommand{\DIFdeltex}[1]{{\protect\color{red}\sout{#1}}}                      %DIF PREAMBLE
%DIF SAFE PREAMBLE %DIF PREAMBLE
\providecommand{\DIFaddbegin}{} %DIF PREAMBLE
\providecommand{\DIFaddend}{} %DIF PREAMBLE
\providecommand{\DIFdelbegin}{} %DIF PREAMBLE
\providecommand{\DIFdelend}{} %DIF PREAMBLE
\providecommand{\DIFmodbegin}{} %DIF PREAMBLE
\providecommand{\DIFmodend}{} %DIF PREAMBLE
%DIF FLOATSAFE PREAMBLE %DIF PREAMBLE
\providecommand{\DIFaddFL}[1]{\DIFadd{#1}} %DIF PREAMBLE
\providecommand{\DIFdelFL}[1]{\DIFdel{#1}} %DIF PREAMBLE
\providecommand{\DIFaddbeginFL}{} %DIF PREAMBLE
\providecommand{\DIFaddendFL}{} %DIF PREAMBLE
\providecommand{\DIFdelbeginFL}{} %DIF PREAMBLE
\providecommand{\DIFdelendFL}{} %DIF PREAMBLE
%DIF HYPERREF PREAMBLE %DIF PREAMBLE
\providecommand{\DIFadd}[1]{\texorpdfstring{\DIFaddtex{#1}}{#1}} %DIF PREAMBLE
\providecommand{\DIFdel}[1]{\texorpdfstring{\DIFdeltex{#1}}{}} %DIF PREAMBLE
\newcommand{\DIFscaledelfig}{0.5}
%DIF HIGHLIGHTGRAPHICS PREAMBLE %DIF PREAMBLE
\RequirePackage{settobox} %DIF PREAMBLE
\RequirePackage{letltxmacro} %DIF PREAMBLE
\newsavebox{\DIFdelgraphicsbox} %DIF PREAMBLE
\newlength{\DIFdelgraphicswidth} %DIF PREAMBLE
\newlength{\DIFdelgraphicsheight} %DIF PREAMBLE
% store original definition of \includegraphics %DIF PREAMBLE
\LetLtxMacro{\DIFOincludegraphics}{\includegraphics} %DIF PREAMBLE
\newcommand{\DIFaddincludegraphics}[2][]{{\color{blue}\fbox{\DIFOincludegraphics[#1]{#2}}}} %DIF PREAMBLE
\newcommand{\DIFdelincludegraphics}[2][]{% %DIF PREAMBLE
\sbox{\DIFdelgraphicsbox}{\DIFOincludegraphics[#1]{#2}}% %DIF PREAMBLE
\settoboxwidth{\DIFdelgraphicswidth}{\DIFdelgraphicsbox} %DIF PREAMBLE
\settoboxtotalheight{\DIFdelgraphicsheight}{\DIFdelgraphicsbox} %DIF PREAMBLE
\scalebox{\DIFscaledelfig}{% %DIF PREAMBLE
\parbox[b]{\DIFdelgraphicswidth}{\usebox{\DIFdelgraphicsbox}\\[-\baselineskip] \rule{\DIFdelgraphicswidth}{0em}}\llap{\resizebox{\DIFdelgraphicswidth}{\DIFdelgraphicsheight}{% %DIF PREAMBLE
\setlength{\unitlength}{\DIFdelgraphicswidth}% %DIF PREAMBLE
\begin{picture}(1,1)% %DIF PREAMBLE
\thicklines\linethickness{2pt} %DIF PREAMBLE
{\color[rgb]{1,0,0}\put(0,0){\framebox(1,1){}}}% %DIF PREAMBLE
{\color[rgb]{1,0,0}\put(0,0){\line( 1,1){1}}}% %DIF PREAMBLE
{\color[rgb]{1,0,0}\put(0,1){\line(1,-1){1}}}% %DIF PREAMBLE
\end{picture}% %DIF PREAMBLE
}\hspace*{3pt}}} %DIF PREAMBLE
} %DIF PREAMBLE
\LetLtxMacro{\DIFOaddbegin}{\DIFaddbegin} %DIF PREAMBLE
\LetLtxMacro{\DIFOaddend}{\DIFaddend} %DIF PREAMBLE
\LetLtxMacro{\DIFOdelbegin}{\DIFdelbegin} %DIF PREAMBLE
\LetLtxMacro{\DIFOdelend}{\DIFdelend} %DIF PREAMBLE
\DeclareRobustCommand{\DIFaddbegin}{\DIFOaddbegin \let\includegraphics\DIFaddincludegraphics} %DIF PREAMBLE
\DeclareRobustCommand{\DIFaddend}{\DIFOaddend \let\includegraphics\DIFOincludegraphics} %DIF PREAMBLE
\DeclareRobustCommand{\DIFdelbegin}{\DIFOdelbegin \let\includegraphics\DIFdelincludegraphics} %DIF PREAMBLE
\DeclareRobustCommand{\DIFdelend}{\DIFOaddend \let\includegraphics\DIFOincludegraphics} %DIF PREAMBLE
\LetLtxMacro{\DIFOaddbeginFL}{\DIFaddbeginFL} %DIF PREAMBLE
\LetLtxMacro{\DIFOaddendFL}{\DIFaddendFL} %DIF PREAMBLE
\LetLtxMacro{\DIFOdelbeginFL}{\DIFdelbeginFL} %DIF PREAMBLE
\LetLtxMacro{\DIFOdelendFL}{\DIFdelendFL} %DIF PREAMBLE
\DeclareRobustCommand{\DIFaddbeginFL}{\DIFOaddbeginFL \let\includegraphics\DIFaddincludegraphics} %DIF PREAMBLE
\DeclareRobustCommand{\DIFaddendFL}{\DIFOaddendFL \let\includegraphics\DIFOincludegraphics} %DIF PREAMBLE
\DeclareRobustCommand{\DIFdelbeginFL}{\DIFOdelbeginFL \let\includegraphics\DIFdelincludegraphics} %DIF PREAMBLE
\DeclareRobustCommand{\DIFdelendFL}{\DIFOaddendFL \let\includegraphics\DIFOincludegraphics} %DIF PREAMBLE
%DIF END PREAMBLE EXTENSION ADDED BY LATEXDIFF

\begin{document}

%DIF <  Interaction of estuaries and fluvial processes in estuarine agricultural areas. 
\DIFdelbegin %DIFDELCMD < 

%DIFDELCMD < %%%
\DIFdelend \section*{Seepage models}

\subsection*{Background}

During inlet closures, the estuary loses water through the wave-built sand berm.

Daily seepage is estimated from a control-volume balance over the basin:
\DIFdelbegin \begin{displaymath}
\DIFdel{S \approx Q_r - \frac{dV}{dt},
}\end{displaymath}%DIFAUXCMD
\DIFdelend \DIFaddbegin \[
\DIFadd{S \approx Q_r - \frac{dV}{dt},
}\]\DIFaddend 
after filtering overwash days and treating evaporation as small relative to $Q_r$ during the analyzed closures. Preliminary analysis indicates seepage is substantial, accounting for roughly 20--60\% of total discharge (report values in cfs).

We use Darcy scaling to represent flow through the berm, treating the berm as a porous body with cross-sectional flow area $A_b$ and flow length $L_b$ between estuary and ocean. The seepage rate $S$ $[L^3T^{-1}]$ scales with the head difference $\Delta h = h_V - h_{\rm ocn}$:
\DIFdelbegin \begin{displaymath}
\DIFdel{S = K\,A_b\,\frac{\Delta h}{L_b},
}\end{displaymath}%DIFAUXCMD
\DIFdelend \DIFaddbegin \[
\DIFadd{S = K\,A_b\,\frac{\Delta h}{L_b},
}\]\DIFaddend 
where $K$ $[LT^{-1}]$ is hydraulic conductivity. Because $A_b$ and $L_b$ can change with water level (e.g., wetted thickness grows while the effective path length/tortuosity evolves), we \DIFdelbegin \DIFdel{absorb }\DIFdelend \DIFaddbegin \DIFadd{group material and }\DIFaddend geometry into an \emph{integrated berm conductance} \DIFdelbegin \begin{displaymath}
\DIFdel{C \;\equiv\; K\,\frac{A_b}{L_b} \qquad }[\DIFdel{L^2T^{-1}}]\DIFdel{,
}\end{displaymath}%DIFAUXCMD
\DIFdel{and model seepage as
}\begin{displaymath}
\DIFdel{S = C(\cdot)\,\Delta h,
}\end{displaymath}%DIFAUXCMD
\DIFdelend \DIFaddbegin \DIFadd{and write
}\[
\DIFadd{\mathcal{C}(\cdot) \;\equiv\; K\,\frac{A_b}{L_b}\qquad }[\DIFadd{L^2T^{-1}}]\DIFadd{,
\qquad
S \;=\; \mathcal{C}(\cdot)\,\Delta h,
}\]\DIFaddend 
where \DIFdelbegin \DIFdel{$C(\cdot)$ }\DIFdelend \DIFaddbegin \DIFadd{the conductance function $\mathcal{C}(\cdot)$ }\DIFaddend may vary with estuary level $h_V$ or with $\Delta h$.

\textbf{Transmissivity connection.} If the wetted thickness is $b$ and the alongshore width of the active flow section is $W_b$, then $A_b \approx b\,W_b$ and
\DIFdelbegin \begin{displaymath}
\DIFdel{C = K\,\frac{A_b}{L_b}
   = (K\,b)\,\frac{W_b}{L_b}
   = T\,\Big(\frac{W_b}{L_b}\Big),
}\end{displaymath}%DIFAUXCMD
\DIFdelend \DIFaddbegin \[
\DIFadd{\mathcal{C}(\cdot) \;=\; K\,\frac{A_b}{L_b}
   \;=\; (K\,b)\,\frac{W_b}{L_b}
   \;=\; T\,\Big(\frac{W_b}{L_b}\Big),
}\]\DIFaddend 
where $T=K\,b$ $[L^2T^{-1}]$ is \DIFdelbegin %DIFDELCMD < {%%%
\DIFdel{transmissivity}%DIFDELCMD < }%%%
\DIFdel{. Thus $C$ is a }\DIFdelend \DIFaddbegin \emph{\DIFadd{transmissivity}}\DIFadd{. Thus $\mathcal{C}$ is }\DIFaddend transmissivity-like multiplied by a geometric \emph{shape factor} $W_b/L_b$. When $W_b/L_b$ is nearly constant, \DIFdelbegin \DIFdel{$C\propto T$}\DIFdelend \DIFaddbegin \DIFadd{$\mathcal{C}\propto T$}\DIFaddend ; when geometry varies with level, the factor $f(\cdot)$ below captures that variation.

We parameterize
\DIFdelbegin \begin{displaymath}
\DIFdel{C(\cdot) = C_{\rm ref}\, f(\cdot),
}\end{displaymath}%DIFAUXCMD
\DIFdelend \DIFaddbegin \[
\DIFadd{\mathcal{C}(\cdot) \;=\; C_{\rm ref}\, f(\cdot),
}\]\DIFaddend 
with $C_{\rm ref}$ $[L^2T^{-1}]$ a baseline (reference) conductance and $f(\cdot)$ dimensionless; $f=1$ in the linear model.

\subsection*{Three candidate models}

\paragraph{(1) Constant \DIFdelbegin \DIFdel{$C$ }\DIFdelend \DIFaddbegin \DIFadd{conductance }\DIFaddend (linear\DIFdelbegin \DIFdel{fit}\DIFdelend )\DIFdelbegin \DIFdel{: constant conductance}\DIFdelend .}
\DIFdelbegin \begin{displaymath}
\DIFdel{f(\cdot)=1
\quad\Rightarrow\quad
S = C_L\,\Delta h,
}\end{displaymath}%DIFAUXCMD
\DIFdelend \DIFaddbegin \[
\DIFadd{f(\cdot)=1
\quad\Rightarrow\quad
S \;=\; C_L\,\Delta h,
}\]\DIFaddend 
with $C_L$ $[L^2T^{-1}]$ constant.

\noindent\textit{Advantage:} simple, transparent, convenient for operations. \\
\textit{Assumption:} the effective seepage cross-section/path is approximately constant over the analyzed range of $\Delta h$.

\paragraph{(2) Visitor-level response (\DIFdelbegin \DIFdel{shifted–power }\DIFdelend \DIFaddbegin \DIFadd{shifted--power }\DIFaddend on $h_V$).}
\DIFdelbegin \begin{displaymath}
\DIFdel{C(\cdot)= C_V\,z_{V}^{\,b_{V}}
\quad\Rightarrow\quad
S = C_V\,\Delta h\, z_{V}^{\,b_{V}},
}\end{displaymath}%DIFAUXCMD
\DIFdelend \DIFaddbegin \[
\DIFadd{\mathcal{C}(\cdot)= C_h\,f_h(h_V),
\qquad
f_h(h_V) \;=\; z_{h}^{\,b_{h}},
}\]\DIFaddend 
\DIFdelbegin \begin{displaymath}
\DIFdel{z_{V}=\frac{h_{V}+\theta_{V}}{\,m_{V}\,},\qquad
m_{V}=\operatorname{median}(h_{V})+\theta_{V}.
}\end{displaymath}%DIFAUXCMD
\DIFdelend \DIFaddbegin \[
\DIFadd{z_{h}\;=\;\frac{h_{V}+\theta_{h}}{\,m_{h}\,},\qquad
m_{h}\;=\;\operatorname{median}(h_{V})+\theta_{h},
}\]\DIFaddend 
\DIFdelbegin \DIFdel{Here $C_V$ }\DIFdelend \DIFaddbegin \[
\DIFadd{\Rightarrow\qquad
S \;=\; C_h\,\Delta h\, z_{h}^{\,b_{h}}.
}\]
\DIFadd{Here $C_h$ }\DIFaddend $[L^2T^{-1}]$ is the conductance at the reference state \DIFdelbegin \DIFdel{$(z_V=1)$; $\theta_V$ }\DIFdelend \DIFaddbegin \DIFadd{$(z_h=1)$; $\theta_h$ }\DIFaddend is an offset (threshold) chosen so \DIFdelbegin \DIFdel{$h_V+\theta_V>0$; $b_V\ge 0$ }\DIFdelend \DIFaddbegin \DIFadd{$h_V+\theta_h>0$; $b_h\ge 0$ }\DIFaddend governs sensitivity to $h_V$.

\paragraph{(3) Head-difference response (\DIFdelbegin \DIFdel{shifted–power }\DIFdelend \DIFaddbegin \DIFadd{shifted--power }\DIFaddend on $\Delta h$).}
\DIFdelbegin \begin{displaymath}
\DIFdel{C(\cdot)= C_\Delta\,z_{\Delta}^{\,b_{\Delta}}
\quad\Rightarrow\quad
S = C_\Delta\,\Delta h\, z_{\Delta}^{\,b_{\Delta}},
}\end{displaymath}%DIFAUXCMD
\DIFdelend \DIFaddbegin \[
\DIFadd{\mathcal{C}(\cdot)= C_\Delta\,f_\Delta(\Delta h),
\qquad
f_\Delta(\Delta h) \;=\; z_{\Delta}^{\,b_{\Delta}},
}\]\DIFaddend 
\DIFdelbegin \begin{displaymath}
\DIFdel{z_{\Delta}=\frac{\Delta h+\theta_{\Delta}}{\,m_{\Delta}\,},\qquad 
m_{\Delta}=\operatorname{median}(\Delta h)+\theta_{\Delta}.
}\end{displaymath}%DIFAUXCMD
\DIFdelend \DIFaddbegin \[
\DIFadd{z_{\Delta}\;=\;\frac{\Delta h+\theta_{\Delta}}{\,m_{\Delta}\,},\qquad 
m_{\Delta}\;=\;\operatorname{median}(\Delta h)+\theta_{\Delta},
}\]\DIFaddend 
\DIFaddbegin \[
\DIFadd{\Rightarrow\qquad
S \;=\; C_\Delta\,\Delta h\, z_{\Delta}^{\,b_{\Delta}}.
}\]
\DIFaddend Here $C_\Delta$ $[L^2T^{-1}]$ is the reference conductance; $\theta_\Delta$ ensures $\Delta h+\theta_\Delta>0$; $b_\Delta\ge 0$ is estimated separately from \DIFdelbegin \DIFdel{$b_V$}\DIFdelend \DIFaddbegin \DIFadd{$b_h$}\DIFaddend .

\subsection*{Bounds in estimating model parameters}

We do \emph{not} force $A_b=0$ at $h_V=0$ (the visitor gauge datum is NAVD88, not estuary depth; $h_V=0$ $\approx$ mean low tide can still have nonzero wetted thickness). Practically, we bound offsets globally so the powered bases remain positive and physically reasonable (e.g., $\theta \ge -\min(x)+\varepsilon$, plus a small site-scale upper cap).

\subsection*{Consistency notes}

\begin{itemize}
\item Use $L_b$ consistently for flow length.
\item \DIFdelbegin \DIFdel{Keep $C$ as conductance $[L^2T^{-1}]$ and $C_{\rm ref}$ as the baseline parameter}\DIFdelend \DIFaddbegin \DIFadd{Use $\mathcal{C}(\cdot)$ for the conductance }\emph{\DIFadd{function}} \DIFadd{and $C_{\square}\in\{C_L,C_h,C_\Delta\}$ for the dimensionful reference conductances $[L^2T^{-1}]$}\DIFaddend .
\item $z$ factors are dimensionless (median \DIFdelbegin \DIFdel{+ }\DIFdelend \DIFaddbegin \DIFadd{$+$ }\DIFaddend offset), so \DIFdelbegin \DIFdel{$C_{\rm ref}$ }\DIFdelend \DIFaddbegin \DIFadd{$C_{\square}$ }\DIFaddend retains physical units.
\end{itemize}
\DIFdelbegin %DIFDELCMD < 

%DIFDELCMD < %%%
\DIFdel{Since river discharge, estuary height and hypsometry (stage-volume relationship) data are available, seepage can be estimated by mass balance. 
}%DIFDELCMD < 

%DIFDELCMD < %%%
\DIFdel{ICEs function as tidal systems when there is an ocean connection, and as salt-stratified lakes during closures,  which prevents vertical mixing and can lead to hypoxia. 
}%DIFDELCMD < 

%DIFDELCMD < %%%
\DIFdel{An integrated hydraulic conductance can be obtained by regression of seepage on the head difference between estuary and ocean.
}%DIFDELCMD < 

%DIFDELCMD < %%%
\DIFdel{Bar-built estuaries with relatively small (cross sectional areab100 m2) and shallow tidal inlets are wide- spread in Mediterranean climates and along wave-exposed coasts.
Small, sandy coastal lagoons are common along the U.S. Pacific Coast, provide valuable habitat for salmonids, and are being squeezed by ongoing sea level rise and coastal development.
}\DIFdelend 

\end{document}
